%
% LISTBOXES
%

% #1 Box Name
% #2 Entry Name
% For use in sentence
\newcommand{\toolaef}[2]{\hyperref[entry:#1-#2]{$#1_{\ref{entry:#1-#2}}$}}
\newcommand{\itemref}[2]{entry \hyperref[entry:#1-#2]{$#1_{\ref{entry:#1-#2}}$}}
% For use inside math
\WithSuffix\newcommand\itemref*[2]{\hyperref[entry:#1-#2]{#1_{\ref{entry:#1-#2}}}}
% For use at start of sentence
\newcommand{\Itemref}[2]{Entry \hyperref[entry:#1-#2]{$#1_{\ref{entry:#1-#2}}$}}

\newcommand{\rboxref}[1]{\hyperref[lb:#1]{#1}}


% #1 - Box Name
% #2 - Entry Name
\newcommand{\entryLabel}[2]{entry:#1-#2}

% #1 - Entry Name
% #2 - Entry Code
\newcommand{\entry}[2]{
    \edef\lastentryid{\theentryid}
    \refstepcounter{entryid}
    \label{\entryLabel{\boxName}{#1}}

    \node[below=1mm of \lastentryid, inner sep=2mm, text width=\boxWidth, rounded corners, thick, draw=blue!20, fill=blue!10] (\theentryid) {\small\begin{obeylines}#2\end{obeylines}};
    \node[below left=0.5mm and 0mm of \theentryid.north east, @link] {\footnotesize $\boxName_\theentryid$};
}

\newcounter{boxId}
\setcounter{boxId}{0}
\newcounter{entryid}
\counterwithin*{entryid}{boxId}

% #1 - Width
% #2- Box Name (used in entry label)
\newenvironment{listbox}[2][4cm]{
    \edef\boxName{#2}
    \edef\boxWidth{#1}
    \refstepcounter{boxId}\label{lb:\boxName}

    \begin{center}
    \begin{tikzpicture}
    \node[color=blue!40] (0) {\small Listbox \boxName};
}{
    \begin{pgfonlayer}{background}
        \draw[draw=blue!30, fill=blue!1, rounded corners, thick]
        ($(current bounding box.south west) - (4pt, 4pt)$)
        rectangle
        ($(current bounding box.north east) + (4pt, 2pt)$);
    \end{pgfonlayer}
    \end{tikzpicture}
    \end{center}
}
